\chapter{Vui mà thấm}
\markboth{Vui mà thấm}{Funny but Lasting}

\section{Đi học mà không mang não thì cặp nặng vô ích.}
\begin{truyen}{Đi học mà không mang não thì cặp nặng vô ích.}{Humorous but deep.}
\chuhoa{C}{ả lớp cười ồ, nhưng thầy muốn nhắc rằng mang đủ sách vở chưa đủ. Điều cần có mặt thật sự trong lớp là sự suy nghĩ của chính mình.}
\textit{(The class burst out laughing, but the teacher wanted to remind them that bringing books is not enough. What truly needs to be present in class is their own thinking.)}
\end{truyen}

\begin{baihoc}
Học không chỉ cần dụng cụ; học cần đầu óc tỉnh táo và tham gia thật.
\textit{(Learning needs more than tools; it needs an awake and engaged mind.)}
\end{baihoc}

\begin{ghinhoanh}
\textbf{Short English to remember:}

Bring your mind to class, not just your bag.

\end{ghinhoanh}

\begin{tuhocnhanh}
1. Em có mang sự tập trung vào lớp không? \textit{Do you bring focus into class?}

2. Điều gì làm em học mà như đang vắng mặt? \textit{What makes you study as if you were absent?}
\end{tuhocnhanh}
\ngancach

\section{Bút em viết được tương lai, đừng chỉ dùng nó để gạch chơi.}
\begin{truyen}{Bút em viết được tương lai, đừng chỉ dùng nó để gạch chơi.}{Humorous but deep.}
\chuhoa{T}{hầy nói vậy khi thấy vài em nghịch bút giữa giờ. Câu đùa khiến lớp cười, nhưng cũng nhắc rằng những công cụ nhỏ hằng ngày có thể trở thành phương tiện để xây một cuộc đời nghiêm túc hơn.}
\textit{(The teacher said this when he saw students doodling during class. The joke made the class laugh, but it also reminded them that small daily tools can become instruments for building a more serious life.)}
\end{truyen}

\begin{baihoc}
Điều nhỏ trong tay em cũng có thể tạo nên điều lớn nếu em dùng đúng.
\textit{(Even a small thing in your hand can create something great if used well.)}
\end{baihoc}

\begin{ghinhoanh}
\textbf{Short English to remember:}

Use your tools to build, not waste.

\end{ghinhoanh}

\begin{tuhocnhanh}
1. Em đang dùng thời gian học để xây hay để giết thời gian? \textit{Are you using study time to build or to kill time?}

2. Việc nhỏ nào em có thể dùng nghiêm túc hơn? \textit{Which small thing can you use more seriously?}
\end{tuhocnhanh}
\ngancach

\section{Học dở không đáng sợ, giả vờ hiểu mới đáng sợ.}
\begin{truyen}{Học dở không đáng sợ, giả vờ hiểu mới đáng sợ.}{Humorous but deep.}
\chuhoa{C}{âu nói làm nhiều em cười gượng. Không hiểu thật thì còn học được; nhưng giả vờ hiểu khiến mình ngừng hỏi, ngừng sửa, và ngừng tiến bộ.}
\textit{(The sentence made many students laugh awkwardly. Not truly understanding is still fixable; pretending to understand makes a person stop asking, stop correcting, and stop growing.)}
\end{truyen}

\begin{baihoc}
Sự thành thật với cái chưa biết luôn tốt hơn cái vỏ bề ngoài thông minh.
\textit{(Honesty about what you do not know is always better than the appearance of being smart.)}
\end{baihoc}

\begin{ghinhoanh}
\textbf{Short English to remember:}

Pretending blocks learning.

\end{ghinhoanh}

\begin{tuhocnhanh}
1. Em có hay gật đầu dù chưa hiểu không? \textit{Do you often nod even when you do not understand?}

2. Vì sao giả vờ hiểu lại nguy hiểm? \textit{Why is pretending to understand dangerous?}
\end{tuhocnhanh}
\ngancach

\section{Mạng yếu còn kết nối lại được, ý chí yếu thì phải tự nạp.}
\begin{truyen}{Mạng yếu còn kết nối lại được, ý chí yếu thì phải tự nạp.}{Humorous but deep.}
\chuhoa{C}{ả lớp phì cười vì câu ví von quá thời đại. Nhưng điều thầy muốn nói rất rõ: khi ý chí xuống thấp, không ai nạp hộ mình được hoàn toàn; mình phải tự dựng lại nhịp sống và động lực.}
\textit{(The class laughed because the comparison felt very modern. But the teacher's point was clear: when your will becomes weak, no one can fully recharge it for you; you must rebuild your rhythm and motivation yourself.)}
\end{truyen}

\begin{baihoc}
Có lúc cần được giúp, nhưng cũng có lúc em phải tự làm mình mạnh lên.
\textit{(Sometimes you need help, but sometimes you must strengthen yourself.)}
\end{baihoc}

\begin{ghinhoanh}
\textbf{Short English to remember:}

Recharge your will, not just your devices.

\end{ghinhoanh}

\begin{tuhocnhanh}
1. Điều gì đang làm ý chí em yếu đi? \textit{What is draining your will?}

2. Em có thể “nạp lại” bản thân bằng cách nào? \textit{How can you recharge yourself?}
\end{tuhocnhanh}
\ngancach

\section{Em không thiếu thời gian, em đang tiêu nhầm chỗ thôi.}
\begin{truyen}{Em không thiếu thời gian, em đang tiêu nhầm chỗ thôi.}{Humorous but deep.}
\chuhoa{N}{ghe xong, nhiều em cười rồi im. Câu nói chạm đúng một điều khó chịu: không phải ngày của mình ngắn hơn người khác, mà là mình để những việc ít giá trị lấy mất quá nhiều giờ.}
\textit{(After hearing it, many students laughed and then went quiet. The sentence touched something uncomfortable: the day is not shorter for you than for others; you simply let low-value things take too many hours.)}
\end{truyen}

\begin{baihoc}
Quản thời gian trước hết là quản sự ưu tiên.
\textit{(Managing time begins with managing priorities.)}
\end{baihoc}

\begin{ghinhoanh}
\textbf{Short English to remember:}

You don't always need more time, just better use of it.

\end{ghinhoanh}

\begin{tuhocnhanh}
1. Thời gian của em đang thất thoát ở đâu? \textit{Where is your time leaking away?}

2. Em sẽ cắt bớt điều gì trước? \textit{What will you reduce first?}
\end{tuhocnhanh}
\ngancach

\section{Thầy không đòi em hoàn hảo; thầy đòi em đừng lười hợp lý hóa.}
\begin{truyen}{Thầy không đòi em hoàn hảo; thầy đòi em đừng lười hợp lý hóa.}{Humorous but deep.}
\chuhoa{C}{âu đùa nghe sắc mà đúng. Nhiều học sinh không lười vì thiếu khả năng, mà vì giỏi tìm lý do khiến sự lười của mình trông có vẻ hợp lý.}
\textit{(The joke sounded sharp but true. Many students are not lazy because they lack ability, but because they are good at inventing reasons that make laziness look reasonable.)}
\end{truyen}

\begin{baihoc}
Thành thật với mình là bước đầu để bỏ thói trì trệ.
\textit{(Being honest with yourself is the first step to leaving stagnation behind.)}
\end{baihoc}

\begin{ghinhoanh}
\textbf{Short English to remember:}

Stop explaining laziness as if it were wisdom.

\end{ghinhoanh}

\begin{tuhocnhanh}
1. Em có hay tự bào chữa cho việc mình chưa cố gắng không? \textit{Do you often justify not trying enough?}

2. Một lý do nào em cần bỏ ngay? \textit{Which excuse should you drop immediately?}
\end{tuhocnhanh}
\ngancach

\section{Cứ bảo “mai học” thì mai sẽ thành một nghề.}
\begin{truyen}{Cứ bảo “mai học” thì mai sẽ thành một nghề.}{Humorous but deep.}
\chuhoa{C}{ả lớp cười vì thấy mình đâu đó trong câu này. Trì hoãn lâu ngày không còn là một hành vi đơn lẻ, nó thành nếp sống và dần dần nuốt mất tương lai.}
\textit{(The class laughed because they recognized themselves in the sentence. Long-term postponement is no longer a single action; it becomes a way of living and slowly eats the future.)}
\end{truyen}

\begin{baihoc}
Đừng để “ngày mai” trở thành nơi em gửi hết trách nhiệm của hôm nay.
\textit{(Do not let “tomorrow” become the place where you dump all of today's responsibility.)}
\end{baihoc}

\begin{ghinhoanh}
\textbf{Short English to remember:}

If you keep saying tomorrow, tomorrow becomes your habit.

\end{ghinhoanh}

\begin{tuhocnhanh}
1. Việc nào em đã hẹn “mai” quá nhiều lần? \textit{Which task have you postponed to “tomorrow” too many times?}

2. Em có thể làm một phần của nó ngay bây giờ không? \textit{Can you do one part of it right now?}
\end{tuhocnhanh}
\ngancach

\section{Bàn học không cắn em đâu, cứ lại gần đi.}
\begin{truyen}{Bàn học không cắn em đâu, cứ lại gần đi.}{Humorous but deep.}
\chuhoa{T}{hầy nói vậy khi thấy có em quanh quẩn mãi mà không ngồi vào bàn. Cả lớp cười, nhưng ai cũng hiểu vấn đề thật là sự ngại bắt đầu chứ không phải cái bàn.}
\textit{(The teacher said this when he saw students hovering around without sitting down to work. The class laughed, but everyone understood that the real problem was hesitation to begin, not the desk itself.)}
\end{truyen}

\begin{baihoc}
Bắt đầu thường khó hơn làm tiếp, nên việc đầu tiên là cứ ngồi xuống đã.
\textit{(Starting is often harder than continuing, so the first task is simply to sit down.)}
\end{baihoc}

\begin{ghinhoanh}
\textbf{Short English to remember:}

Start first; courage often follows action.

\end{ghinhoanh}

\begin{tuhocnhanh}
1. Điều gì làm em ngại ngồi vào bàn học? \textit{What makes you avoid your study desk?}

2. Em có thể giảm lực cản bắt đầu thế nào? \textit{How can you reduce the resistance to starting?}
\end{tuhocnhanh}
\ngancach

\section{Chưa hiểu mà gật đầu nhiều quá cũng không thành thông minh.}
\begin{truyen}{Chưa hiểu mà gật đầu nhiều quá cũng không thành thông minh.}{Humorous but deep.}
\chuhoa{C}{âu này khiến lớp cười nhưng cũng chạm vào thói quen rất phổ biến: gật cho qua, gật cho xong, rồi về nhà mới thấy lỗ hổng to đùng.}
\textit{(The line made the class laugh but also touched a common habit: nodding to get through the moment, then going home and finding a huge gap in understanding.)}
\end{truyen}

\begin{baihoc}
Đừng dùng cái gật đầu để che đi điều em còn chưa rõ.
\textit{(Do not use a nod to hide what you still do not understand.)}
\end{baihoc}

\begin{ghinhoanh}
\textbf{Short English to remember:}

Nodding is not understanding.

\end{ghinhoanh}

\begin{tuhocnhanh}
1. Em có dám dừng lại để hỏi cho rõ không? \textit{Do you dare to stop and ask for clarity?}

2. Một cái gật đầu giả tạo gây hại thế nào? \textit{How can a fake nod do harm?}
\end{tuhocnhanh}
\ngancach

\section{Muốn điểm cao mà học kiểu đoán lòng đề thì hơi mạo hiểm.}
\begin{truyen}{Muốn điểm cao mà học kiểu đoán lòng đề thì hơi mạo hiểm.}{Humorous but deep.}
\chuhoa{C}{ả lớp bật cười vì ai cũng từng hy vọng đề ra trúng phần mình học lướt. Nhưng học kiểu cầu may không phải chiến lược bền; nó chỉ làm mình lo lắng hơn trước mỗi kỳ kiểm tra.}
\textit{(The class laughed because everyone had once hoped the exam would match the part they skimmed. But studying by luck is not a lasting strategy; it only creates more anxiety before every test.)}
\end{truyen}

\begin{baihoc}
Học chắc luôn yên tâm hơn học kiểu đánh cược.
\textit{(Solid learning is always calmer than gambling on luck.)}
\end{baihoc}

\begin{ghinhoanh}
\textbf{Short English to remember:}

Luck is not a study plan.

\end{ghinhoanh}

\begin{tuhocnhanh}
1. Em đang học chắc hay học may rủi? \textit{Are you studying solidly or by chance?}

2. Điều gì giúp em bớt phụ thuộc vào may mắn? \textit{What helps you rely less on luck?}
\end{tuhocnhanh}
\ngancach